\section{Introduction} \label{intro}
The foundation of QKD traces back to Stephen Wiesner's (then unpublished) \textit{Conjugate Coding} manuscript where it was proposed that encoding information in incompatible bases can be used for secure communications. He noted that wave guides could be used as a communication channel to send quantum information encoded in the polarization of light. Messages can be constructed using either the vertical-horizontal or circular polarization basis. Consider the usual computational (Z) basis: $\left\{\ket{0}, \ket{1}\right\}$. The two basis states can be represented by photons polarized in the horizontal or vertical direction. These polarizations are orthogonal, so a polarimeter can be oriented such that photons from one of the two states will be detected with certainty. Now, consider the X basis: $\left\{\ket{-}, \ket{+}\right\}$. The two basis states can be represented by photons polarized in the left-handed circular or right-handed circular direction. Again, these two polarizations are orthogonal, so measuring in the X basis means a set of photons polarized in this basis can be effectively filtered. However, if the measurement device is set to measure in the X basis while the photons are prepared in the Z basis, or vice-versa, the probability of the measurement device detecting any of the prepared photons is 50 (e.g. \% $\abs{\bra{0}\ket{+}}^2 = \flatfrac{1}{2}$). Wiesner noted that the two bases are non-commuting observables akin to position and momentum. Only measuring the received information using the correct basis can a deterministic result be discerned. By measuring in the wrong basis, not only is the outcome probabilistic, but uncertainty is introduced because the original state is disturbed by the measurement and no further knowledge about it can be recovered, since identical copies of the unknown, original state can't be created (in accordance with the no-cloning theorem). Additionally, the fact that a state can't be recovered makes the detection of eavesdropping possible by the receiver. Hence, the probabilistic nature of incompatible measurement bases provides a foundation for key distribution protocols because they provide a means of encoding and verifying information. The principles outlined in this manuscript lie at the heart of QKD protocols later developed. % Cite Wiesner's manuscript for this paragraph

The first QDK protocol was developed in 1984 by Bennett and Brassard. Traditional public key cryptosystems require having a pair of keys, one public key that can be shared over a communication channel, and one private key that must be kept secret for the message to remain secure. The algorithm proposed by Bennett and Brassard creates a shared private key between the sender and receiver, akin to private key cryptosystems. However, unlike private key cryptosystems, the BB84 protocol maintains one of the advantages of public key cryptosystems, which is the distribution of secure random keys over an insecure (quantum) communications channel without the need of a secret key shared a priori, as long as an authenticated classical communication channel is available.



